% 建模活动单 · 第九章
% 主题：期望与决策——离散概率模型。
% 两个活动分别用离散分布（抽奖）与二项分布（猜题）走“建模 → 计算 → 检验 → 决策”的过程。
% 文中数值均经 Python 实算核验。
\section{活动十八：抽奖游戏设计——期望与方差}

\subsection{情境}
某商场举办抽奖活动：共发行 1000 张奖券，其中头奖 1 张、奖金 1000 元；
二奖 9 张、各 100 元；三奖 90 张、各 10 元；其余 900 张无奖。
每张奖券售价 2 元。

\subsection{提出问题}
\begin{enumerate}
  \item 买一张奖券，期望能得多少奖金？期望净赚多少？
  \item 这个游戏“值得参与”吗？
  \item 商场（庄家）的期望利润是多少？这门生意做得下去吗？
\end{enumerate}

\subsection{选模型}
把“一张奖券的奖金”看成离散型随机变量 $X$，先列出分布列：
\begin{center}
\footnotesize
\setlength{\tabcolsep}{4pt}
\renewcommand{\arraystretch}{1.4}
\begin{tabular}{p{2.6cm}p{2.6cm}p{3.2cm}}
\toprule
奖金 $X$（元） & 张数 & 概率 $P(X)$ \\
\midrule
$1000$ & $1$ & $0.001$ \\
$100$ & $9$ & $0.009$ \\
$10$ & $90$ & $0.090$ \\
$0$ & $900$ & $0.900$ \\
\bottomrule
\end{tabular}
\end{center}
离散型随机变量的期望与方差：
\[
  E(X)=\sum_i x_ip_i,\qquad D(X)=E(X^2)-[E(X)]^2.
\]

\subsection{计算}
期望奖金：
\[
  E(X)=1000\times0.001+100\times0.009+10\times0.090+0\times0.900=2.8\text{（元）}.
\]
二阶矩：
\[
  E(X^2)=1000^2\times0.001+100^2\times0.009+10^2\times0.090=1000+90+9=1099,
\]
\[
  D(X)=1099-2.8^2=1099-7.84=1091.16,\qquad \sigma=\sqrt{1091.16}\approx33.03\text{（元）}.
\]
设净收益 $Y=X-2$（元），则
\[
  E(Y)=2.8-2=0.8\text{（元）},\qquad D(Y)=D(X)=1091.16
\]
（减去常数不改变方差）。

庄家一侧：卖出 1000 张，收入 $1000\times2=2000$ 元；期望派奖
$1000\times2.8=2800$ 元，故期望利润
\[
  2000-2800=-800\text{（元）},
\]
按收入计的利润率为 $-800/2000=-40\%$，返奖率高达 $140\%$。

\subsection{讨论：值得参与吗}
\begin{itemize}
  \item 对玩家：期望净收益 $+0.8$ 元看似划算，但 $90\%$ 的奖券一无所获，
        中头奖的概率仅 $0.001$；标准差 $33.03$ 元约是期望净收益的 $41$ 倍。
        “平均赚 $0.8$ 元”是大量重复之后的结果，只买一张几乎必然亏钱。
  \item 对庄家：每场期望亏 800 元，这个定价显然不可持续。保本票价是 $2.8$ 元；
        若想获得 $10\%$ 的利润率，票价应定为 $2.8/0.9\approx3.11$ 元。
        例如定价 $3$ 元时，玩家期望净收益 $-0.2$ 元，庄家利润率
        $200/3000\approx6.7\%$，就成了一门“正常”的生意。
  \item 所以“值得参与吗”没有唯一答案：它取决于你的立场（玩家还是庄家）、
        购买次数和能承受多大的波动。数学期望只给平均值，不给保证。
\end{itemize}
\key{期望描述大量重复的平均结果，方差描述波动大小；期望为正，不等于单次不亏。}

\section{活动十九：考试猜题策略——二项分布与正态近似}

\subsection{情境}
某次测验共 20 道四选一单选题，每题 5 分，满分 100 分，通过线为 60 分。
一名考生对所有题目都没有把握，每题只能随机猜测。

\subsection{提出问题}
\begin{enumerate}
  \item 全部随机猜，期望得多少分？
  \item 通过（不少于 60 分）的可能性有多大？
  \item 用正态分布近似二项分布来估算，可靠吗？
\end{enumerate}

\subsection{选模型}
每题独立，答对概率 $p=\dfrac14$。设答对题数为 $X$，则
\[
  X\sim B\!\left(20,\frac14\right),\qquad
  E(X)=np=5,\qquad D(X)=npq=20\times\frac14\times\frac34=3.75,
\]
标准差 $\sigma=\sqrt{3.75}\approx1.94$。总得分 $S=5X$；
通过当且仅当 $S\ge60$，即 $X\ge12$（至少答对 12 题）。

\subsection{计算}
期望得分
\[
  E(S)=5E(X)=25\text{（分）},
\]
距离通过线 60 分相差悬殊——平均只能拿到四分之一的分。

精确通过概率（二项分布）：
\[
  P(X\ge12)=\sum_{k=12}^{20}\binom{20}{k}\left(\frac14\right)^k\left(\frac34\right)^{20-k}\approx0.00094,
\]
即约 $0.094\%$，不到千分之一。

正态近似：常用判据是 $np\ge5$ 且 $nq\ge5$。这里 $np=5$、$nq=15$，
形式上恰好踩在门槛上。使用连续性修正：
\[
  z=\frac{11.5-5}{\sqrt{3.75}}\approx3.36,\qquad
  P(X\ge12)\approx1-\Phi(3.36)\approx0.00040.
\]
两种算法结论的方向一致（通过希望极其渺茫），但数值相差约 $2.4$ 倍：
正态近似把精确值 $0.00094$ 低估成了 $0.00040$。
原因是 $n=20$ 太小，而且我们关心的是远离均值 $5$ 的右尾，
那里的近似误差被明显放大。

\subsection{检验与反思}
再看一个对照：若通过线降为 40 分（需答对 8 题），精确值
$P(X\ge8)\approx0.1018$；正态近似给出
$z=(7.5-5)/\sqrt{3.75}\approx1.29$，$1-\Phi(1.29)\approx0.0984$，两者就很接近。
可见近似的可靠性取决于所求位置离均值有多远：靠近均值时可用，深入尾部时
应回到精确计算。

进一步看，“蒙”永远不是策略。把每题正确率提高后，通过概率的变化如下：
\begin{center}
\footnotesize
\setlength{\tabcolsep}{4pt}
\renewcommand{\arraystretch}{1.4}
\begin{tabular}{p{2.6cm}p{3.0cm}p{3.4cm}}
\toprule
每题正确率 $p$ & 期望得分 & 通过概率 $P(X\ge12)$ \\
\midrule
$\dfrac14$（全猜） & $25$ & $0.00094$ \\
$\dfrac12$ & $50$ & $0.252$ \\
$\dfrac23$ & $66.7$ & $0.809$ \\
\bottomrule
\end{tabular}
\end{center}
真正决定通过率的是每题的正确率，而不是猜测技巧；
“期望”和“尾部概率”要放在一起看，才能做出靠谱的决定。
\key{用正态近似前先查 $np,nq$；还要检查所求概率是否落在远离均值的尾部。}
